\chapter{THIẾT KẾ HỆ THỐNG}

\section{Giới thiệu}

Sau khi xác định đầy đủ các yêu cầu trong tài liệu SRS, chương này sẽ trình bày toàn bộ quá trình thiết kế hệ thống Chatbot hỗ trợ học tập ứng dụng mô hình Retrieval-Augmented Generation (RAG). Nội dung cốt lõi bao gồm thiết kế kiến trúc phần mềm, xây dựng biểu đồ lớp (Class Diagram), thiết kế mô hình thực thể liên kết (ERD) cho cơ sở dữ liệu và đặc tả luồng dữ liệu (Data Flow) giữa các thành phần.

Việc thiết kế hệ thống được định hướng theo tiêu chí đảm bảo khả năng mở rộng linh hoạt, dễ dàng bảo trì và tối ưu cho việc tích hợp thêm các mô hình ngôn ngữ hoặc kỹ thuật Embedding mới. Đồng thời, kiến trúc phải hỗ trợ tốt cho module Benchmark để phục vụ mục tiêu đánh giá và so sánh các cấu hình RAG khác nhau.

\section{Kiến trúc tổng thể}

Hệ thống được xây dựng theo mô hình Client-Server hiện đại, kết hợp chặt chẽ với kiến trúc AI Pipeline chuyên biệt cho RAG. Các thành phần chính cấu thành nên hệ thống bao gồm:
\begin{itemize}
	\item \textbf{Frontend ReactJS}: Chịu trách nhiệm render giao diện tương tác, biểu đồ Dashboard và xử lý logic phía người dùng.
	\item \textbf{Backend FastAPI}: Đóng vai trò là Controller trung tâm, nhận request, điều phối các tác vụ AI và giao tiếp với các Database.
	\item \textbf{ChromaDB}: Cơ sở dữ liệu Vector chuyên dụng để lưu trữ và truy xuất các Embedding sinh ra từ tài liệu.
	\item \textbf{LLM (Gemini/OpenAI)}: Đóng vai trò sinh ngôn ngữ, nhận Prompt từ hệ thống và tạo câu trả lời cuối cùng.
	\item \textbf{Module Benchmark}: Thực thi các kịch bản kiểm thử, tính toán điểm số (Accuracy, Relevancy) và lưu trữ kết quả cấu hình tối ưu.
	\item \textbf{Cơ sở dữ liệu MySQL}: Lưu trữ siêu dữ liệu (metadata), thông tin người dùng, lịch sử đoạn chat và các báo cáo phân tích.
\end{itemize}

\section{Thiết kế lớp (Class Diagram)}

Biểu đồ lớp (Class Diagram) mô tả cấu trúc hướng đối tượng của toàn bộ Backend hệ thống, thể hiện các đối tượng tham gia vào quy trình xử lý và mối quan hệ giữa chúng.

\begin{figure}[H]
	\centering
	\includegraphics[width=\textwidth]{images/class_diagram.png}
	\caption{Biểu đồ lớp của hệ thống}
	\label{fig:class}
\end{figure}

Hệ thống được thiết kế xung quanh các lớp chính yếu sau:

\subsection*{SinhVien \& QuanTriVien}
Hai lớp này kế thừa từ một lớp User trừu tượng chung, đại diện cho những người sử dụng hệ thống. Lớp `SinhVien` cung cấp các phương thức để gửi câu hỏi, nhận câu trả lời và tra cứu lịch sử hội thoại. Lớp `QuanTriVien` được cấp quyền cao hơn, bổ sung các phương thức để quản lý danh sách tài liệu môn học, theo dõi biểu đồ thống kê hệ thống và khởi chạy các kịch bản kiểm thử quy mô lớn.

\subsection*{HeThongChatbot}
Đây là lớp điều phối trung tâm của toàn bộ ứng dụng (Facade/Controller). Lớp này nhận đầu vào là câu hỏi từ `CauHoi`, sau đó tự động điều hướng luồng xử lý: gọi module xử lý văn bản, kết nối đến `MoHinhAI` để tạo Embedding, truy vấn `TaiLieuMonHoc` (thông qua ChromaDB), và cuối cùng là sinh Prompt để gọi LLM trả về kết quả.

\subsection*{TaiLieuMonHoc}
Lớp này đóng gói các thuộc tính và hành vi liên quan đến một tài liệu tải lên (PDF, DOCX). Nó đảm nhiệm việc đọc luồng dữ liệu file, thực hiện chia nhỏ văn bản (Chunking) theo các chiến lược được thiết lập sẵn, và giao tiếp với cơ sở dữ liệu Vector để lưu trữ các đoạn text đã được phân mảnh.

\subsection*{MoHinhAI}
Đây là một Interface trừu tượng đại diện cho các dịch vụ AI bên ngoài (như Gemini, GPT, hoặc các mô hình Embedding mã nguồn mở). Lớp này quy định chuẩn giao tiếp chung cho việc gửi Prompt và nhận Text, hoặc gửi Text và nhận Vector, giúp hệ thống dễ dàng thay thế hoặc nâng cấp mô hình AI trong tương lai mà không cần sửa đổi các lớp khác.

\subsection*{KetQuaSoSanh \& BaoCaoPhanTich}
Các lớp này thuộc về module Benchmark. `KetQuaSoSanh` thực hiện tính toán độ đo (Metrics) giữa các cấu hình Chunking và Embedding khác nhau dựa trên tập test set. `BaoCaoPhanTich` sau đó sẽ tổng hợp các số liệu này (Accuracy, Latency, Relevancy) để kết xuất thành báo cáo và cung cấp dữ liệu cho việc vẽ biểu đồ Radar trên Frontend.

\section{Thiết kế cơ sở dữ liệu}

Hệ thống kết hợp sử dụng cơ sở dữ liệu quan hệ (MySQL) cho dữ liệu có cấu trúc và cơ sở dữ liệu Vector (ChromaDB) cho dữ liệu Embedding. Biểu đồ ERD dưới đây mô tả cấu trúc của hệ cơ sở dữ liệu quan hệ.

\begin{figure}[H]
	\centering
	\includegraphics[width=\textwidth]{images/erd.png}
	\caption{ERD của hệ thống}
	\label{fig:erd}
\end{figure}

Các bảng dữ liệu cốt lõi bao gồm:

\subsection*{Users, Questions, Answers}
Bảng `Users` lưu trữ thông tin xác thực và phân quyền của sinh viên và quản trị viên. Bảng `Questions` lưu lại lịch sử câu hỏi do người dùng gửi lên, liên kết 1-n với bảng `Users`. Bảng `Answers` lưu trữ câu trả lời tương ứng được AI sinh ra, liên kết 1-1 với `Questions`. Cấu trúc này cho phép khôi phục toàn bộ bối cảnh của một phiên trò chuyện (Session).

\subsection*{Documents \& Document Chunks}
Bảng `Documents` quản lý metadata của tài liệu gốc (tên file, đường dẫn lưu trữ, ngày tải lên, môn học). Bảng `Document Chunks` (được đồng bộ logic với ChromaDB) lưu thông tin về các đoạn văn bản đã bị cắt nhỏ từ `Documents`. Bảng này giúp hệ thống đối chiếu và truy xuất ngược lại vị trí gốc của đoạn văn (ví dụ: nằm ở trang số mấy của file PDF nào) để tạo chức năng trích dẫn nguồn cho Chatbot.

\subsection*{Chunking Strategies \& Embedding Models}
Đây là các bảng cấu hình tham số. `Chunking Strategies` lưu trữ các quy tắc chia văn bản (ví dụ: Recursive, Fixed-size, Semantic). `Embedding Models` lưu trữ danh sách các mô hình tạo vector được hệ thống hỗ trợ. Hai bảng này là cơ sở để người dùng lựa chọn cấu hình trong giao diện Chat Sandbox.

\subsection*{Experiments, Evaluations \& Reports}
Cụm bảng này phục vụ cho tính năng Benchmark. `Experiments` đại diện cho một lần chạy thử nghiệm với một cấu hình (Chunking + Embedding) cụ thể. `Evaluations` lưu điểm số chi tiết của từng tiêu chí (Accuracy, Faithfulness, Latency) cho từng câu hỏi trong tập test. Cuối cùng, `Reports` sẽ lưu kết quả thống kê tổng hợp (Average Score) của toàn bộ đợt thử nghiệm đó để hiển thị lên Dashboard.

\section{Luồng xử lý dữ liệu (RAG Pipeline)}

Mọi truy vấn của người dùng đều trải qua một quy trình xử lý khép kín và tự động. Dưới đây là mô tả chi tiết luồng dữ liệu chuẩn (Standard RAG Pipeline) của hệ thống:

\begin{enumerate}
	\item \textbf{Tiếp nhận}: Frontend ReactJS gửi câu hỏi của sinh viên kèm theo cấu hình RAG hiện tại (nếu có) thông qua REST API tới Backend FastAPI.
	\item \textbf{Tạo Vector câu hỏi}: Backend sử dụng `MoHinhAI` (được chỉ định trong cấu hình) để chuyển đổi nội dung câu hỏi thành một Vector đa chiều.
	\item \textbf{Truy xuất (Retrieval)}: Backend kết nối tới ChromaDB, thực thi thuật toán Cosine Similarity để so sánh vector câu hỏi với hàng ngàn vector tài liệu có sẵn, từ đó rút trích ra Top-K đoạn văn bản (chunks) có ý nghĩa gần gũi nhất.
	\item \textbf{Xây dựng Prompt}: Các đoạn văn bản Top-K này được Backend tổng hợp, làm sạch và chèn vào một Template Prompt chuẩn, yêu cầu AI chỉ được phép trả lời dựa trên lượng thông tin này.
	\item \textbf{Sinh kết quả (Generation)}: Prompt hoàn chỉnh được gửi qua API tới mô hình LLM (ví dụ Gemini 2.5 Flash). LLM suy luận và sinh ra câu trả lời cuối cùng.
	\item \textbf{Phản hồi}: Backend ghi nhận lịch sử vào cơ sở dữ liệu MySQL, tính toán thời gian phản hồi (Latency), gán thông tin nguồn trích dẫn và trả về cho ReactJS để render ra giao diện.
\end{enumerate}

\section{Kết luận chương}

Chương này đã phác thảo bức tranh kỹ thuật toàn diện của hệ thống thông qua việc trình bày kiến trúc Client-Server kết hợp AI Pipeline, mô hình hướng đối tượng Class Diagram và cơ sở dữ liệu quan hệ ERD. Việc thiết kế lấy RAG và Benchmark làm trung tâm không chỉ đảm bảo ứng dụng chạy mượt mà, mà còn tạo nền tảng vững chắc để đo lường, tinh chỉnh và tối ưu hóa hệ thống ở giai đoạn triển khai kế tiếp.